%% Nessa semana vamos continuar nosso projeto. Vocês lembraram de fazer as reuniões diárias semana passada?
%% Na Sprint planning, quem ficou responsável por fazer o quê? Vocês verificaram o andamento das tarefas nas reuniões diárias?
%% - O que vocês já tem pronto? Nessa semana, vocês já devem ter pronta ao menos a história e a definição do que deve aparecer quadro a quadro.
%% É importante que nessa semana vocês também façam testes das gerações dos quadrinhos. Vocês vão reparar que tem um mundo de possibilidade, ele é super interessante, mas que a IA tem limitações de quantos quadros pode-se gerar por dia (sim, por isso HQ são tão boas para a gente).
%% Nos veremos semana que vem com esses quadrinhos bem encaminhados.
%% Ahh, e atenção: Se a sua atividade não for feita em grupo, você não vai disputar as pontuações completas dos eventos do Sprint, somente das tarefas.

\documentclass[11pt]{article}
\usepackage[a4paper,includeheadfoot,margin=2.54cm]{geometry}
\usepackage{parskip}

% fonte adotada
\usepackage{mathpazo}
\usepackage{tgpagella}
%\usepackage{domitian}
\usepackage[T1]{fontenc}
\let\oldstylenums\oldstyle

\usepackage{parskip}
\usepackage{graphicx}
\usepackage{enumitem}
\usepackage{pifont}
\usepackage{tikz}
\usetikzlibrary{calc, shapes}
\usepackage{amssymb}
\usepackage{amsmath}
\usepackage{minted}
\usepackage{hyperref} % cria os links de referencia das equacoes

\renewcommand{\labelenumii}{\arabic{enumi}.\arabic{enumii}} % enumeracao em subniveis

\begin{document}

\underline{\textbf{Ata - Reuniões Diárias}}\par
\textbf{Universidade Aberta Capixaba - UNAC/ Instituto Federal do Espírito Santo - IFES}\\
\textbf{Engenharia de Software}\par
Prof. Formador Carlos Lins Borges Azevedo\\

Anderson A. Fraga (20252TSIS0258)\\
João Vitor N. Carlini ()\\

\begin{enumerate}
    \item Ata 1 - Reunião 25/04/26
        \begin{itemize}
            \item Descrição: Delimitação do que foi realizado na Sprint inicial e organização do que faremos na atual Sprint (\textit{sprint retrospective}).
            \item Registro da reunião:
                \begin{table}[!h]
                    \hline
                    registro 1& registro 2 & registro 3 \\
                    \hline
                \end{table}
            \item Próximos passos: 
        \end{itemize}
\end{enumerate}

\end{document}
